\documentclass[11pt]{article}
% Shared preamble for all daily exercises
% Update this file to change settings (padding, parindent, etc.) for all exercises

\usepackage[margin=0.75in]{geometry}
\usepackage{amsmath}
\usepackage{amssymb}

% Custom settings
\setlength{\parindent}{0pt}
\setlength{\parskip}{0.2cm}

\title{Daily Exercise}
\author{Dhruman Gupta}
\date{20th March}

\begin{document}

% Common header for daily exercises
% This file can be included in other LaTeX documents

\maketitle

\noindent
\textbf{Course:} CS-3330, Theory of Computation

\vspace{0.2cm}

\noindent
\textbf{Question:}\noindent 
Prove that Context Free Languages are not closed under quotienting.

\textbf{Answer:}

Take right quotient to mean
\[
L_1 / L_2 = \{x \mid \exists y \in L_2 \text{ such that } xy \in L_1\}.
\]

We will give a context free language $P$ and a language $K$ such that $P/K$ is not context free.

Let $\Sigma = \{a,b,c\}$ and let $\#$ be a symbol not in $\Sigma$. Define
\[
P = \{w\#w^R \mid w \in \Sigma^*\}.
\]
This language is context free, since it is generated by the grammar
\[
S \to aSa \mid bSb \mid cSc \mid \#.
\]

Now let
\[
A = \{a^n b^n c^n \mid n \geq 0\}.
\]
We know that $A$ is not context free. Define
\[
K = \{\#w^R \mid w \in A\}.
\]

We claim that $P/K = A$.

Take any $x \in \Sigma^*$. Then
\[
x \in P/K \iff \exists y \in K \text{ such that } xy \in P.
\]
Since every $y \in K$ has the form $\#w^R$ for some $w \in A$, this becomes
\[
x \in P/K \iff \exists w \in A \text{ such that } x\#w^R \in P.
\]
But $P$ consists exactly of strings of the form $u\#u^R$. Therefore,
\[
x\#w^R \in P \iff x = w.
\]
Hence
\[
x \in P/K \iff x \in A.
\]
So $P/K = A$.

Since $A$ is not context free, $P/K$ is not context free. But $P$ is context free. Therefore context free languages are not closed under quotienting.

\end{document}
